% Supplementary material of the AnytimeSwarm-CC manuscript (material moved out of the main text when it was shortened).
% Build after main: pdflatex supplementary && bibtex supplementary && pdflatex supplementary && pdflatex supplementary
% Every number comes from the generated tables/numbers*.tex files; none is typed by hand.
\documentclass[preprint,11pt]{elsarticle}
\usepackage[T1]{fontenc}
\usepackage{newtxtext}
\usepackage[varvw]{newtxmath}
\usepackage{graphicx}
\usepackage{booktabs}
\usepackage{amsmath}
\usepackage{xcolor}
\usepackage{xr-hyper}
\usepackage[hidelinks]{hyperref}
\usepackage[margin=2.5cm]{geometry}
\externaldocument{main}
\biboptions{sort&compress}
\newtheorem{prop}{Proposition}
\renewcommand{\theprop}{S\arabic{prop}}
\newproof{proof}{Proof}
\newcommand{\todo}[1]{\textcolor{red}{[TODO: #1]}}
\IfFileExists{tables/static_numbers.tex}{\input{tables/static_numbers.tex}}{}
\IfFileExists{tables/static_numbers_XL.tex}{\input{tables/static_numbers_XL.tex}}{}
\IfFileExists{tables/static_numbers_scale.tex}{\input{tables/static_numbers_scale.tex}}{}
\IfFileExists{tables/numbers.tex}{\input{tables/numbers.tex}}{}
\IfFileExists{tables/numbers_s2.tex}{\input{tables/numbers_s2.tex}}{}
\IfFileExists{tables/numbers_s3.tex}{\input{tables/numbers_s3.tex}}{}
\IfFileExists{tables/numbers_s3x.tex}{\input{tables/numbers_s3x.tex}}{}
\IfFileExists{tables/numbers_err.tex}{\input{tables/numbers_err.tex}}{}
\IfFileExists{tables/numbers_s4.tex}{\input{tables/numbers_s4.tex}}{}
\IfFileExists{tables/numbers_s4x.tex}{\input{tables/numbers_s4x.tex}}{}
\IfFileExists{tables/numbers_s4c.tex}{\input{tables/numbers_s4c.tex}}{}
\newcommand{\val}[1]{\ifcsname v#1\endcsname\csname v#1\endcsname\else\todo{#1}\fi}
\newcommand{\fig}[2]{\IfFileExists{figures/#1.pdf}{\includegraphics[width=#2]{figures/#1.pdf}}{\fbox{\todo{figure \detokenize{#1}}}}}
\newcommand{\inputtable}[1]{\IfFileExists{tables/#1.tex}{\input{tables/#1.tex}}{\todo{table \detokenize{#1}}}}
\renewcommand{\thesection}{S\arabic{section}}
\renewcommand{\thetable}{S\arabic{table}}
\renewcommand{\thefigure}{S\arabic{figure}}
\renewcommand{\theequation}{S\arabic{equation}}
\journal{Robotics and Autonomous Systems}

\begin{document}
\begin{frontmatter}
\title{Supplementary material for: Lead-time-scaled chance constraints with a nominal fallback for fast
multi-UAV replanning under uncertain execution}
\author{Cao-Phuc Ha, Anh Nguyen Vu Minh, Phu-Nguyen Le}
\begin{abstract}
This document contains the parts of the evaluation that the main paper only summarises: the stopping rule and its two
delay costs with their properties, the scenario generator and the fitted error models, the exact-execution benchmark of
the stopping rules, the full results of study~1, further results of study~2 per regime and in tier C, and the
reliability of the error model in study~4. Section, equation and table numbers without the prefix S refer to the main paper.
\end{abstract}
\end{frontmatter}

\section{Stopping rule and delay costs}\label{sec:Sstop}
\subsection{Learned stopping rule}
Every 5\,ms, the stopping rule compares the predicted gain of searching for another $H=10$\,ms with the cost of doing so.
With $\widehat{\Delta J}$ as the predicted gain, the search stops at time $\tau$ if
\begin{equation}\label{eq:Sstop}
\widehat{\Delta J}\;\le\;C_{\mathrm{delay}}(\tau)+\mu W H .
\end{equation}
The gain is predicted from the search history by a Poisson generalised linear model with a log
link~\cite{mccullagh1989glm} and a small $L_2$ penalty. Its eleven features describe the elapsed time, the number of
iterations, the recent improvements over 5\,ms, over 20\,ms and over the second half of the elapsed time, the time since
the last improvement, the total improvement over the repaired plan, the unassigned share of the reward, the numbers of
tasks and UAVs, and the iteration rate. The target is the improvement in the next 10\,ms in percent of $W$. The model is
fitted on 500\,ms search trajectories of development scenarios only, and a prediction costs one dot product. The search
loop enforces the hard time budget $B$ independently of the stopping rule. The price of computation is
$\mu=10^{-5}$ per millisecond.

\subsection{Deterministic delay cost}
The delay cost is computed from the best plan found so far. A late start of a route is partly absorbed by the waiting
times along it, in the same way as the forward time slack of routing with time windows~\cite{savelsbergh1992vrptw}. The
tolerance $\mathrm{tol}_j$ of a task is the margin between its finish time and its deadline, plus all waiting times on
the route up to this task. A task is lost when the start shift $\sigma_j(\tau)$ of its route exceeds its tolerance. The
start shift is zero for a UAV that is still busy, and it equals the waiting time until the publication for an idle UAV.
With $v_j=\alpha w_j-\delta\,\mathbb{1}[\text{$j$ changes UAV}]$ as the value of task $j$ in the objective and
$\Lambda(\tau)=\sum_j v_j\,\mathbb{1}[\sigma_j(\tau)>\mathrm{tol}_j]$, the delay cost of searching $H$ milliseconds longer
is $C_{\mathrm{delay}}(\tau)=\Lambda(\tau+H)-\Lambda(\tau)$.

\begin{prop}\label{prop:delay}
For a feasible plan executed exactly, (i) $\Lambda(\tau)$ is the value lost through missed deadlines if every UAV
follows its route, (ii) the loss of the objective is at most $\Lambda(\tau)$ if late tasks are skipped, and
(iii) $\Lambda(\tau)$ does not decrease with $\tau$ and is zero while all UAVs are busy with their committed tasks.
Hence $C_{\mathrm{delay}}$ is non-negative and vanishes while no UAV is waiting.
\end{prop}
\begin{proof}
A route that starts later by some shift reaches its first task later by the same amount. Each waiting time absorbs part
of the delay, so the delay at the end of a task equals the start shift minus all earlier waiting times, if this
difference is positive. The task therefore misses its deadline exactly when the start shift exceeds its tolerance.
Waiting consumes no energy, which proves (i). If late tasks are skipped, the triangle inequality ensures that every
later task is reached no later and with no more energy. All tasks that are on time in (i) remain on time, and the flown
distance, hence the energy term $\gamma E$, does not grow. The objective therefore loses the values $v_j$ of the skipped
tasks and nothing else, which is at most $\Lambda(\tau)$ and proves (ii). Part (iii) holds because the start shift does
not decrease with $\tau$ and every tolerance is non-negative in a feasible plan.
\end{proof}

\subsection{Execution-aware delay cost}\label{sec:Srisk}
Under uncertain execution, task $j$ is lost when $\sigma_j(\tau)+e_j>\mathrm{tol}_j$, where $e_j$ is its finish-time
error. The execution-aware cost replaces the indicator in $\Lambda$ by the probability of this event. The error of task
$j$ is modelled as $e_j\sim\mathcal{N}(m(h_j),s(h_j)^2)$, where $h_j=f_j-t_0$ is its lead time, the nominal finish time
minus the time of the planning snapshot. The mean and the spread are linear in the lead time,
\begin{equation}\label{eq:Serrmodel}
m(h)=b_0+b_1h,\qquad s(h)=s_0+s_1h,
\end{equation}
because errors accumulate along a route. The expected dissemination latency $\bar L$ is added to the start shift. The
expected loss and the execution-aware delay cost are
\begin{equation}\label{eq:Srisk}
\Lambda_{\mathrm{R}}(\tau)=\sum_j v_j\,\Phi\!\left(\frac{\sigma_j(\tau)+\bar L+m(h_j)-\mathrm{tol}_j}{s(h_j)}\right),
\qquad C_{\mathrm{R}}(\tau)=\Lambda_{\mathrm{R}}(\tau+H)-\Lambda_{\mathrm{R}}(\tau),
\end{equation}
where $\Phi$ is the standard normal distribution function. Search and predictor are unchanged, and $C_{\mathrm{R}}$
replaces $C_{\mathrm{delay}}$ in Eq.~\eqref{eq:Sstop}.

\begin{prop}\label{prop:risk}
(i) With $b_0=b_1=s_0=s_1=0$ and $\bar L=0$, $\Lambda_{\mathrm{R}}=\Lambda$, and the stopping rule takes exactly the
decisions of the deterministic one. (ii) $\Lambda_{\mathrm{R}}(\tau)$ does not decrease with $\tau$, so $C_{\mathrm{R}}\ge0$.
(iii) Under the error model, $\Lambda_{\mathrm{R}}(\tau)$ is the expected value of the tasks whose deadline is missed if
every UAV follows its route.
\end{prop}
\begin{proof}
With zero spread, $\Phi$ is read as the indicator of a positive argument, which gives (i). The start shift does not
decrease with $\tau$ and $\Phi$ is non-decreasing, which gives (ii). Part (iii) follows from the proof of
Proposition~\ref{prop:delay}, because task $j$ misses its deadline exactly when its delay at the finish,
$\sigma_j+e_j$ minus the absorbed waiting time, is positive, and the expectation of the indicator of this event is the
probability in Eq.~\eqref{eq:Srisk}.
\end{proof}

The parameters of Eq.~\eqref{eq:Serrmodel} are fitted on development missions. The mean is estimated by least squares,
and the spread by least squares on the absolute residuals scaled by $\sqrt{\pi/2}$, which is less sensitive to the heavy
tail of the errors than squared residuals. The latency $\bar L$ is the mean of the latency distribution in tier B and
the measured mean plan-acceptance latency in tier C.



\section{Scenario generator, error models and exact-execution snapshots}\label{sec:Sscen}
\subsection{Scenarios}
Missions are created by a seeded generator with the parameters in Table~\ref{tab:params}. The map is a square with a
side of 10\,km and the depot at its centre, and half of the tasks are clustered around three hotspots. Half of the tasks
are known at the start, and the others arrive in batches. Three fleet scales with 5, 10 and 20 UAVs and 25, 50 and 140
tasks are combined with three event profiles. The arrivals profile contains six arrival batches, the failures profile
three batches and two UAV failures, and the deadlines profile three batches and four tightened deadlines. A fourth scale
with 40 UAVs and 280 tasks is used for scalability only.

\begin{table}[t]
\caption{Scenario generator parameters (distances in m, times in s, energy in metres of flight).}\label{tab:params}
\centering\small
\setlength{\tabcolsep}{4pt}
\begin{tabular}{ll@{\qquad}ll}
\toprule
Map side / horizon & 10\,000 / 1\,800 & Service time & $U[30,90]$ \\
UAV speed & $U[10,16]$\,m/s & Time-window slack & $U[180,600]$ \\
Energy capacity & $U[18,28]$\,km & Release delay & $U[0,120]$ \\
Energy reserve & 10\% of capacity & Rewards 1, 2, 3, 5 & probabilities .4, .3, .2, .1 \\
Hover cost & 6 per second & Hotspots & 3, std.\ 600, 50\% of tasks \\
Tasks known at start & 50\% & Scales (UAVs/tasks) & 5/25, 10/50, 20/140, 40/280 \\
\bottomrule
\end{tabular}
\end{table}

\subsection{Static snapshots with exact execution}
The stopping rules are first compared without execution uncertainty. Each mission is simulated once with a
deterministic reference planner, and every replanning call after the start is stored with its previous plan. All
methods are evaluated on these snapshots, and regret is measured against the best known objective of each snapshot,
which includes a 3\,s ALNS run and, for the small and medium scales, the OR-Tools routing solver~\cite{ortools} with
guided local search~\cite{voudouris1999gls}. Because all ALNS variants follow the same search trajectory, whole families
of each rule are evaluated exactly offline, with 12 fixed budgets, 10 stagnation patience values and 9 compute prices,
and each baseline is compared with AnytimeSwarm at the same mean computation. This benchmark uses its own pre-registered
test seeds 2000--2029 (\val{NsnapTest} snapshots of 270 missions), and it was reproduced on the final software platform
before the execution-aware study (Section~\ref{sec:repro}).



\inputtable{tab_errmodels}

\section{Exact execution: stopping rules at equal computation}\label{sec:Sstatic}
\label{sec:static}
On the \val{NsnapTest} test snapshots, regret-2 repair finished within about one millisecond but left a regret of
\val{RepairRegret}\%. With a budget of 250\,ms, ALNS with the full budget reached \val{FixedRegret}\%, whereas
AnytimeSwarm reached \val{ASRegret}\% with \val{ASUsed}\% of the budget, that is, with \val{SavedPct}\% less computation.
Because the stopping rules use different amounts of computation, they are best compared at equal computation.
Table~\ref{tab:scale} makes this comparison for each fleet scale, with the fixed budget and the stagnation rule evaluated
at the mean computation that AnytimeSwarm used on the same scale. AnytimeSwarm reached a lower regret on every scale,
and all paired intervals lie below zero. The reduction was \val{ScMediumRedMin}\% to \val{ScMediumRedMax}\% on the
medium scale, \val{ScLargeRedMin}\% to \val{ScLargeRedMax}\% on the large scale and \val{ScXlargeRedMin}\% to
\val{ScXlargeRedMax}\% with 40 UAVs, a scale that was never used for training. Read in the other direction, the two
baselines needed between \val{ScMediumSpeedMin} and \val{ScSmallSpeedMax} times the computation of AnytimeSwarm to
reach its regret. Over all scales, AnytimeSwarm reached \val{EqRegretAS}\% against \val{EqRegretBaseMin}\% to
\val{EqRegretBaseMax}\% for four baselines that include a fixed budget and a stagnation patience tuned per fleet scale
on development data.

\inputtable{tab_static_scale}

Fig.~\ref{fig:static} shows where this advantage comes from. On the medium and large scales, the quality--compute
frontier of AnytimeSwarm lies below those of fixed budgets and stagnation stopping over the range of practical
computation (Fig.~\ref{fig:static}a). Grouped by the improvement that a 500\,ms search can still find
(Fig.~\ref{fig:static}b), AnytimeSwarm spent \val{AllocASNone}\,ms on snapshots that cannot be improved and
\val{AllocASHigh}\,ms on those with the largest available improvement. Only \val{AllocASShare}\% of its computation went
to snapshots without any possible gain, compared with more than half for the baselines. The learned predictor is
essential for this allocation. It reached an area under the ROC curve~\cite{fawcett2006roc} of \val{AUCLearned} for the
event that the next 10\,ms bring an improvement, compared with \val{AUCOnline} for an extrapolation of the recent
improvement rate.

\begin{figure}[t]
\centering
\begin{minipage}[b]{0.62\linewidth}\centering\fig{fig_static_frontier}{\linewidth}\\[-1pt]{\small (a)}\end{minipage}\hfill
\begin{minipage}[b]{0.36\linewidth}\centering\fig{fig_static_allocation}{\linewidth}\\[-1pt]{\small (b)}\end{minipage}
\caption{Exact execution, test seeds 2000--2029: (a) quality--compute frontiers on identical search trajectories (open
markers: baselines tuned per scale; star: selected AnytimeSwarm setting); (b) mean computation per snapshot, grouped by
the improvement that a 500\,ms search can still find.}\label{fig:static}
\end{figure}

\subsection{Reproduction on the evaluation platform}\label{sec:repro}
The exact-execution benchmark was produced on Windows. Before the execution-aware study, its development pipeline was
re-run on the platform of all later experiments (Ubuntu 24.04 under WSL2, identical Python and package versions). The
deterministic part was reproduced bit for bit, except for differences of one unit in the last place of some energy
values, which stem from the floating-point libraries of the two operating systems. The search ran
\val{JReproSpeedup}\% more iterations per millisecond on the new platform. The improvement predictor, re-trained on the
new timing traces, reached the same compute price by the pre-registered rule, and the conclusions of the
equal-computation comparison agreed in sign and significance in \val{JReproMatched} comparisons. The frozen predictor
was used unchanged in all later experiments.


\section{Study 1: execution-aware stopping on the nominal planning model}\label{sec:Sstudy1}
Study~1 compares AnytimeSwarm-R with the stopping rules of the first group of Section~\ref{sec:methods} on the nominal
planning model, on held-out seeds 3000--3029 in tier B (270 missions over three scales and three profiles, and 90
missions with 40 UAVs) and 4000--4009 in tier C (30 missions per method). Its primary matrix crosses two execution
levels (low, mid), two latency levels (none, high), two latency factors ($\kappa_B=10, 50$) and two triggers (event,
periodic every 60\,s). Its primary hypothesis is that AnytimeSwarm-R serves more reward than ALNS with 40\,ms,
AnytimeSwarm and the stagnation rule.
\subsection{Tier B primary matrix}\label{sec:tierb}
The primary matrix comprises \val{JNRunsMain} mission runs, that is, \val{JNMain} held-out missions in
\val{JNConfigsMain} configurations with eight methods. Table~\ref{tab:main} lists utility, lost tasks and computation per
method, and Fig.~\ref{fig:main} relates utility to computation and shows the difference in every configuration.

Search pays off clearly under uncertain execution. Regret-2 repair reached a utility of \val{JURepair}\%, and
AnytimeSwarm-R reached \val{JURisk}\%, which is \val{JDmRepair} percentage points (pp) or \val{JURelRepair}\% more. The
full budget of 250\,ms was the weakest search method, with \val{JUFixed}\%, and AnytimeSwarm-R served \val{JDFixed}\,pp
more reward than it. The loss of the long budget stems from late plans. A task counts as lost to delay when it is part of
the computed plan but can no longer be served on time when the plan is released. With $\kappa_B=50$, the full budget
lost \val{JLostKFiftyFixed} tasks per mission in this way, against \val{JLostKFiftyRisk} for AnytimeSwarm-R, that is,
\val{JLostRedFixed}\% fewer. Without its delay cost, AnytimeSwarm lost \val{JLostKFiftyNoDelay} tasks per mission in this
way, but its utility was not lower (\val{JUNoDelay}\% against \val{JUAS}\%). The delay cost thus prevents losses from
late plans without a measurable gain in utility on this matrix.

Among the remaining search methods, the differences are small. All seven search methods lie within
\val{JUSpreadSearch}\,pp of each other. AnytimeSwarm-R reached the highest mean utility with the least computation of the adaptive rules,
\val{JMsRisk}\,ms per call against \val{JMsAS}\,ms for AnytimeSwarm and \val{JMsStag}\,ms for the stagnation rule. The
reduction of \val{JMsRedAS}\% relative to AnytimeSwarm is significant (\val{JDMsAS}\,ms per call), so the second
hypothesis holds. The primary hypothesis does not hold at the pre-registered level. The pooled advantage of
AnytimeSwarm-R was \val{JDFixedForty}\,pp over ALNS with 40\,ms, \val{JDAS}\,pp over AnytimeSwarm and \val{JDStag}\,pp
over the stagnation rule, with Holm-adjusted $p$-values of \val{JPFixedForty}, \val{JPAS} and \val{JPStag}. Only the
first interval excludes zero before correction. Per configuration and without correction for multiple testing,
AnytimeSwarm-R was better than ALNS with 40\,ms in \val{JWinFixedForty} of the \val{JNConfigsMain} configurations, better
than the stagnation rule in \val{JWinStag} and never different from AnytimeSwarm, which is consistent with chance
(Fig.~\ref{fig:main}b). The ablation that keeps only the latency term reached \val{JURiskZero}\%, close to AnytimeSwarm.
Together with the misspecification runs below, this suggests that a non-zero error spread mainly makes the rule stop
earlier, rather than that the calibration of the spread matters.

The fixed budget of 40\,ms is a strong reference, because it was matched to the mean computation of AnytimeSwarm on
development data. Its weakness lies in the allocation across fleet scales. AnytimeSwarm-R used \val{JMsScSmallRisk}\,ms
per call on the small scale and \val{JMsScLargeRisk}\,ms on the large scale. On the large scale, it served
\val{JDScLargeFixedForty}\,pp more reward than the budget of 40\,ms and was on a par with the full budget
(\val{JDScLargeFixed}\,pp). On the small scale, it served \val{JDScSmallFixed}\,pp more than the full budget and was on a
par with the budget of 40\,ms (\val{JDScSmallFixedForty}\,pp). A single fixed budget is therefore either too long for
small fleets or too short for large ones, whereas the adaptive rules find a suitable amount of computation on each scale
without tuning. The advantage over ALNS with 40\,ms was \val{JDFFKappa10}\,pp for $\kappa_B=10$ and
\val{JDFFKappa50}\,pp for $\kappa_B=50$, and \val{JDFFExecLow}\,pp under low and \val{JDFFExecMid}\,pp under mid
disturbances.

\inputtable{tab_main}

\begin{figure}[t]
\centering
\fig{fig_main}{\linewidth}
\caption{Tier B primary matrix, held-out seeds 3000--3029: (a) mean utility against mean planning time per call over
all \val{JNConfigsMain} configurations (regret-2 repair, at \val{JURepair}\% and 1\,ms, is not shown); (b) paired
difference of AnytimeSwarm-R minus ALNS with 40\,ms and minus AnytimeSwarm per configuration (execution level, latency
level, $\kappa_B$, trigger) with 95\% bootstrap intervals over \val{JNMain} missions.}\label{fig:main}
\end{figure}

\subsection{Replanning triggers}\label{sec:res-triggers}
The trigger mattered more than the stopping rule. Replanning every 60\,s in addition to the events raised the number of
planning calls from \val{JCallsEvent} to \val{JCallsPeriodic} per mission and the utility of every method, by
\val{JTrigMin}\,pp to \val{JTrigMax}\,pp, with all intervals above zero (\val{JTrigRisk}\,pp for AnytimeSwarm-R,
\val{JTrigFixedForty}\,pp for ALNS with 40\,ms). The third hypothesis therefore holds. Periodic replanning probably helps
because it corrects the accumulated deviations of the executed routes, which the event trigger only notices at the next
event. The advantage of AnytimeSwarm-R over ALNS with 40\,ms was \val{JDFixedFortyEvent}\,pp with the event trigger and
\val{JDFixedFortyPeriodic}\,pp with the periodic trigger. The two values are close, and the protocol did not test their
difference.

\subsection{Scalability to 40 UAVs}\label{sec:res-xl}
With 40 UAVs and 280 tasks, \val{JNXL} held-out missions were run under mid disturbances in \val{JXLNConf}
configurations of latency level and latency factor. All five methods reached between \val{JXLUAS}\% and
\val{JXLUStag}\%. AnytimeSwarm-R reached \val{JXLURisk}\% with \val{JXLMsRisk}\,ms per call. Its pooled difference was
\val{JXLDFixedOneEighty}\,pp to ALNS with 180\,ms, \val{JXLDAS}\,pp to AnytimeSwarm and \val{JXLDStag}\,pp to the
stagnation rule. After Holm correction over the four comparisons of the pre-registered analysis, none of them is
significant ($p_{\mathrm{Holm}}=$ \val{JXLPAS} for AnytimeSwarm), so the fifth hypothesis is not supported either.
AnytimeSwarm-R again needed less computation than AnytimeSwarm (\val{JXLDMsAS}\,ms per call). The stagnation rule
stopped earliest, after \val{JXLMsStag}\,ms, and was significantly better than AnytimeSwarm-R in \val{JXLLossStag} of the
\val{JXLNConf} configurations, which had no latency and $\kappa_B=50$. On this scale, the rules used between \val{JXLMsStag} and 250\,ms
per call, and utility hardly changed within this range of computation. The equal-computation advantage that
AnytimeSwarm showed with 40 UAVs under exact execution (Table~\ref{tab:scale}) therefore did not translate into a
utility gain under uncertain execution.

\subsection{Sensitivity analyses}\label{sec:res-sens}
The sensitivity runs use the held-out missions of the primary matrix without latency and are descriptive. They vary the
execution level beyond the primary matrix, the latency factor and the trigger in \val{JNSensConf} configurations, and
the error model in \val{JNMisConf} further configurations.

Under high disturbances, which lie outside the range of the primary matrix but within that of the fitted error models,
utility fell to about \val{JSensUHighRisk}\% for all methods. The combination of high disturbances and a slow planner
($\kappa_B=50$) gave the largest advantage of the risk-aware cost in this study. AnytimeSwarm-R served
\val{JSensHighKFiftyAS}\,pp more reward than AnytimeSwarm and used \val{JSensMsHighKFiftyRisk}\,ms instead of
\val{JSensMsHighKFiftyAS}\,ms per call. Its difference to ALNS with 40\,ms was \val{JSensHighKFiftyFixedForty}\,pp and to
the stagnation rule \val{JSensHighKFiftyStag}\,pp, so the gain over AnytimeSwarm brought it to the level of the
baselines rather than beyond them. With $\kappa_B=1$, planning time hardly delays the plan, both delay costs are close to
zero because they scale with the shift $\kappa\tau$, and the two rules behave alike. AnytimeSwarm-R then differed from
AnytimeSwarm by \val{JSensKOneAS}\,pp and \val{JSensKOneMsDiff}\,ms per call.

Replanning when a UAV deviated by more than 30\,s from its planned finish time raised the number of calls from
\val{JSensCallsEvent} to \val{JSensCallsDev} per mission and the utility by \val{JSensTrigDevRisk}\,pp for
AnytimeSwarm-R and \val{JSensTrigDevFixedForty}\,pp for ALNS with 40\,ms. Like the periodic trigger, it helped both
stopping rules by a similar amount. Under the deviation trigger, the difference of AnytimeSwarm-R was
\val{JSensDevFixedForty}\,pp to ALNS with 40\,ms and \val{JSensDevAS}\,pp to AnytimeSwarm.

The utility of AnytimeSwarm-R was insensitive to the error model. With the error model of the low, mid or high level
under low and mid disturbances, its difference to AnytimeSwarm ranged between \val{JMisAllLow}\,pp and
\val{JMisAllMid}\,pp, the correctly specified model was not better than a wrong one, and the mean utilities of all
versions lay within \val{JMisSpread}\,pp. Setting the error model to zero gave
\val{JMisAllZero}\,pp. The four misspecification configurations also belong to the primary matrix, so the correctly
specified method was run twice on the same missions. The two runs differ only in the measured planning times, because
the stopping rule reads a wall clock. They gave identical utility in \val{JRepeatSameShare}\% of the missions, and the
mean utility of a method differed by up to \val{JRepeatDiffMax}\,pp between the runs. Differences of this size between
stopping rules should therefore not be read as systematic.

\subsection{Tier C: PX4 software-in-the-loop with Gazebo}\label{sec:res-tierc}
All \val{JNCRuns} held-out flights of tier C were completed, that is, \val{JNC} missions with four methods each. Two runs
did not reach their altitude within 90\,s and were repeated as specified. A further \val{JCNInvalid} runs had a real-time
factor between \val{JCRTFInvalidMin} and \val{JCRTFInvalidMax}, below the validity range of 1.6 to 2.2 stated in the
protocol. The protocol did not specify how to treat such runs, so they were repeated once with the same seed, like a
technical failure, and the original runs were kept. The repeated runs reached the same utility as the original ones, and
the comparisons below are identical with either set (\val{JCDFixedFortyOrig}\,pp and \val{JCDASOrig}\,pp with the
original runs). The real-time factor of the analysed runs was \val{JCRTF} on average and at least \val{JCRTFmin}. A
released plan reached all vehicles after \val{JCLat}\,s of simulation time on average, which is included in the expected
latency of AnytimeSwarm-R.

Table~\ref{tab:tierc} summarises the results. AnytimeSwarm-R reached a utility of \val{JCURisk}\%, against
\val{JCUFixedForty}\% for ALNS with 40\,ms, \val{JCUAS}\% for AnytimeSwarm and \val{JCUFixed}\% for the full budget.
The paired differences were \val{JCDFixedForty}\,pp to ALNS with 40\,ms and \val{JCDAS}\,pp to AnytimeSwarm, with
Holm-adjusted $p$-values of \val{JCPFixedForty} and \val{JCPAS}. The fourth hypothesis is therefore not supported. With
\val{JNC} missions, the intervals are much wider than in tier B and cannot resolve differences of the size seen there.
Two effects of tier B reappear in tier C. The full budget again had the lowest mean utility, although its interval
includes zero (\val{JCDFixed}\,pp to AnytimeSwarm-R), and it was the only method that lost a noticeable number of tasks
to late plans,
\val{JCDropFixed} per mission against at most \val{JCDropAS} for the others. AnytimeSwarm-R again computed least among
the adaptive rules, with \val{JCMsRisk}\,ms per call against \val{JCMsAS}\,ms for AnytimeSwarm, a reduction of
\val{JCMsRedAS}\% (\val{JCDMsAS}\,ms).

\inputtable{tab_tierc}

The same \val{JNC} missions were also run in tier B with the rescaled parameters of tier C. Tier B predicted the outcome
of a flight well. The utility of a run in tier C correlated with that of the same run in tier B with $r=$ \val{JBCr}
(\val{JBCn} runs), and the mean utility per mission with $r=$ \val{JBCrMission} (Fig.~\ref{fig:bc}b). Tier C served
slightly more reward on average, \val{JBCUC}\% against \val{JBCUB}\%, so the kinematic model of tier B was, if anything,
slightly pessimistic for these missions. The small
differences between the methods were not predicted as well. The method means of the two tiers had a Kendall rank
correlation of \val{JBCTau}, and the sign of the per-mission difference between AnytimeSwarm-R and ALNS with 40\,ms
agreed in \val{JBCSignAgreeFixedForty}\% of the missions, counting ties as agreement only when both differences are zero.
Part of the per-run correlation is by construction, because both tiers draw the wind, speed and service disturbances
from the same random numbers, and much of it reflects the different difficulty of the missions. The correlation was
first computed with an erroneous pairing of runs in the analysis script, which was corrected after the results had
been seen; the hypothesis tests were not affected. In tier B, the advantage of AnytimeSwarm-R on these missions
was \val{JBgzDFixedForty}\,pp over ALNS with 40\,ms and \val{JBgzDAS}\,pp over AnytimeSwarm (Fig.~\ref{fig:bc}a). Both
tiers thus agree that the methods differ by at most a few percentage points on this scale, and that the full budget is
the weakest search method, while the order of the three short-computation methods is within the noise of both tiers.

\begin{figure}[t]
\centering
\fig{fig_bc}{\linewidth}
\caption{Tier B against tier C on the same \val{JNC} held-out missions (three UAVs, mid disturbances, $\kappa_B=50$):
(a) paired difference of AnytimeSwarm-R minus each baseline with 95\% bootstrap intervals; (b) utility of every run in
tier C against the same run in tier B (dashed: equality).}\label{fig:bc}
\end{figure}


\section{Study 2: regimes and software-in-the-loop details}\label{sec:Sstudy2}
Fig.~\ref{fig:Ss2main} shows the paired differences of Table~\ref{tab:s2main} of the main paper together with the
computation of every method.

\begin{figure}[t]
\centering
\fig{fig_s2_main}{\linewidth}
\caption{Study~2, tier B, held-out seeds 5000--5029: (a) utility of the proposed replanner minus each method, pooled
over \val{KNConfigs} configurations, with 95\% paired bootstrap intervals over \val{KNMissions} missions; (b) mean
utility against mean computation per call. Ablations are omitted.}\label{fig:Ss2main}
\end{figure}

\subsection{Regimes}\label{sec:s2-regime}
Table~\ref{tab:s2regime} and Fig.~\ref{fig:s2regime} break the utility down by disturbance level, time-window factor and
latency factor. Per configuration and without correction for multiple testing, the proposed replanner was better than
ALNS with 40\,ms in \val{KWinForty} of the \val{KNConfigs} configurations and worse in \val{KLossForty}, better than the
constant buffer in \val{KWinBuf} and worse in \val{KLossBuf}, and better than the conventional planner in \val{KWinConv}
and worse in \val{KLossConv}. The advantage over ALNS with 40\,ms was \val{KDFortyKTen}\,pp, \val{KDFortyKFifty}\,pp and
\val{KDFortyKTwoHundred}\,pp for $\kappa_B=10$, 50 and 200, and the advantage over online prediction stayed between
\val{KDOnlKTen}\,pp and \val{KDOnlKTwoHundred}\,pp.

The advantage was largest with the original windows. Under high disturbances, the proposed replanner served
\val{KDFortyWOneHigh}\,pp more reward than ALNS with 40\,ms and \val{KDBufWOneHigh}\,pp more than the constant buffer.
The losses all occurred under high disturbances with halved windows. In this regime, the proposed replanner reached
\val{KUPropWHalfHigh}\%, \val{KDBufWHalfHigh}\,pp relative to the constant buffer and \val{KDFortyWHalfHigh}\,pp relative
to ALNS with 40\,ms. It served \val{KServedPropWHalfHigh} tasks per mission, of which \val{KLatePropWHalfHigh} late,
against \val{KServedBufWHalfHigh} and \val{KLateBufWHalfHigh} for the constant buffer. The chance-constrained model was
thus too cautious when the windows were short compared with the predicted finish-time error. The effective deadline of
many tasks then probably falls before the earliest possible arrival, so the planner gives up tasks that a less cautious
plan still serves in time. The value of $z$ was selected per disturbance level on development missions, averaged over
both window widths. Selecting it also per window width, or adapting it online, is a natural extension.

\inputtable{tab_s2_regime}

\begin{figure}[t]
\centering
\fig{fig_s2_regime}{\linewidth}
\caption{Study~2, tier B: utility per regime as a function of the latency factor $\kappa_B$, for mid and high
disturbances with the original and with halved time windows (mean over \val{KNMissions} held-out missions per point).
The proposed replanner is shown with stars.}\label{fig:s2regime}
\end{figure}

\subsection{Tier C: PX4 software-in-the-loop with Gazebo}\label{sec:Ss2-tierc}
All \val{KCRuns} held-out flights of study~2 were completed, that is, \val{KCN} missions with four methods each. One run
failed to start and was repeated. As specified in the protocol, \val{KCInvalid} runs with a real-time factor outside
1.6 to 2.2 were repeated once with the same seed, and the original runs were kept. The laptop was also used for other
work during the campaign, and many repeats ran slower than the original runs. Of the \val{KCRuns} analysed runs,
\val{KCValid} lie in the validity range. The real-time factor was \val{KCRTF} on average, with quartiles of
\val{KCRTFq} and a minimum of \val{KCRTFmin}. Two sensitivity analyses that were not pre-registered therefore repeat the
comparisons with the original runs instead of the repeats, and with only those pairs of runs that both lie in the
validity range.

Table~\ref{tab:s2tierc} summarises the results. The proposed replanner reached a utility of \val{KCUProp}\%, against
\val{KCUForty}\% for ALNS with 40\,ms, \val{KCUBuf}\% for the constant buffer and \val{KCUConv}\% for the conventional
planner. The paired differences were \val{KCDForty}\,pp to ALNS with 40\,ms and \val{KCDConv}\,pp to the conventional
planner, with Holm-adjusted $p$-values of \val{KCPForty} and \val{KCPConv}. The fourth hypothesis is therefore not
supported. The secondary comparison with the constant buffer gave \val{KCDBuf}\,pp. The sensitivity analyses lead to
the same conclusion. With the original runs, the differences were \val{KCDFortyOrig}\,pp and \val{KCDConvOrig}\,pp.
With valid runs only, they were \val{KCDFortyValid}\,pp on \val{KCNFortyValid} missions and \val{KCDConvValid}\,pp on
\val{KCNConvValid} missions. No interval in the three analyses excludes zero.

\inputtable{tab_s2_tierc}

The mechanism seen in tier B reappears in the counts of late and skipped tasks. The proposed replanner served
\val{KCLateProp} tasks late per mission, against \val{KCLateForty} for ALNS with 40\,ms, which is \val{KCLateRedForty}\%
fewer. It skipped no task at departure, whereas ALNS with 40\,ms skipped \val{KCSkipForty} and the conventional planner
\val{KCSkipConv} tasks per mission. It computed \val{KCMsProp}\,ms per call, against \val{KCMsForty}\,ms for the fixed
budget. The conventional planner called the search \val{KCCallsConv} times per mission instead of about
\val{KCCallsProp}, as in tier B.


The same \val{KCN} missions were also run in tier B with the rescaled parameters of tier C. As in study~1, tier B
predicted the utility of a flight closely, with a correlation of $r=$ \val{KBCr} over \val{KBCn} runs. On these
missions, however, even tier B did not resolve the differences between the methods (Fig.~\ref{fig:s2bc}). Its estimates
were \val{KBgzDForty}\,pp to ALNS with 40\,ms, \val{KBgzDBuf}\,pp to the constant buffer and \val{KBgzDConv}\,pp to the
conventional planner. The tier C intervals overlap them. The per-mission difference to ALNS with 40\,ms had a standard
deviation of \val{KCSdForty}\,pp in tier C. Confirming the tier B estimate of \val{KCEffForty}\,pp on these missions with
a power of 80\% would require about \val{KCNPowerForty} missions per method. Tier C thus shows that the proposed
replanner runs on an autopilot and a physics engine with the same computation and fewer late tasks, but it cannot
confirm the size of the gain measured in tier B.

\begin{figure}[t]
\centering
\fig{fig_s2_bc}{0.62\linewidth}
\caption{Study~2, tier B against tier C on the same \val{KCN} held-out missions (three UAVs, mid disturbances,
$\kappa_B=50$): paired difference of the proposed replanner minus each baseline with 95\% bootstrap intervals.}
\label{fig:s2bc}
\end{figure}


\section{Study 3: regimes}\label{sec:Sstudy3}
Fig.~\ref{fig:Ss3} shows the paired differences of study~3 per regime.

\begin{figure}[t]
\centering
\fig{fig_s3}{\linewidth}
\caption{Study~3, tier B, held-out seeds 7000--7029: utility of the replanner with nominal fallback minus each other
method, per regime (mean over $\kappa_B$, \val{LNMissions} missions, 95\% paired bootstrap intervals).}\label{fig:Ss3}
\end{figure}

\section{Study 4: reliability of the error model}\label{sec:Sstudy4}
Fig.~\ref{fig:Scoverage} groups the tasks that were served in study~4 by the probability of meeting the deadline that
the error model of their disturbance level predicted at the last published plan, and compares it with the share served
on time. The plans of the proposed model and of the nominal model are evaluated with the same error model. Under mid
disturbances, the share served on time was \val{QCovObsCCMid}\% for the proposed model, against a mean predicted
probability of \val{QCovPredCCMid}\%. Under high disturbances, it was \val{QCovObsCCHigh}\% against
\val{QCovPredCCHigh}\%. Among the tasks with a predicted probability of at least one half, the shares were
\val{QCovObsAboveCCMid}\% and \val{QCovObsAboveCCHigh}\%, against predicted means of \val{QCovPredAboveCCMid}\% and
\val{QCovPredAboveCCHigh}\%. The model was thus slightly optimistic under mid and conservative under high disturbances,
and it ranked the tasks by their risk in both cases.

\begin{figure}[t]
\centering
\fig{fig_coverage}{\linewidth}
\caption{Study~4, tier B, held-out seeds 8000--8029: share of served tasks that met their deadline against the
probability predicted by the error model of the disturbance level, in deciles, for the proposed model with fallback and
the nominal model. The dotted line marks perfect calibration.}\label{fig:Scoverage}
\end{figure}

\bibliographystyle{elsarticle-num}
\bibliography{references}
\end{document}
